\section{Tổng quan lý thuyết}

\noindent Chương này trình bày các nền tảng lý thuyết được sử dụng trong đề tài, bao gồm mô hình học tăng cường, tiến trình quyết định Markov và các thuật toán điều khiển chính sách/giá trị liên quan.

\section{Thuật toán DQN trong học tăng cường}

\noindent Deep Q-Network (DQN) là phương pháp học tăng cường theo giá trị (\textit{dựa trên giá trị}), trong đó tác nhân học hàm giá trị hành động $Q(s, a)$ để ước lượng mức lợi ích kỳ vọng khi chọn hành động $a$ tại trạng thái $s$.

\noindent\textbf{Mục tiêu tối ưu.}

\noindent Với hệ số chiết khấu $\gamma$, mục tiêu của DQN là cực đại hóa tổng phần thưởng tích lũy chiết khấu:
\[
G_t = \sum_{k=0}^{\infty}\gamma^k r_{t+k}
\]
Tác nhân lựa chọn hành động theo chính sách tham lam trên hàm Q:
\[
\pi(s) = \arg\max_{a} Q(s, a)
\]

\noindent\textbf{Xấp xỉ hàm Q bằng mạng neural.}

\noindent Thay vì lưu bảng Q truyền thống (khó mở rộng khi không gian trạng thái lớn), DQN sử dụng mạng neural tham số $\theta$ để xấp xỉ:
\[
Q(s, a; \theta)
\]
Đầu vào là vector trạng thái hệ thống (các đặc trưng từ MDP), đầu ra là giá trị Q tương ứng cho từng hành động rời rạc trong không gian hành động.

\noindent\textbf{Hàm mục tiêu Bellman và hàm mất mát.}

\noindent Tại mỗi bước cập nhật, mục tiêu Bellman được xây dựng:
\[
y = r + \gamma \max_{a'} Q_{\theta^-}(s', a')
\]
trong đó $\theta^-$ là tham số của \textit{mạng mục tiêu}. Hàm mất mát bình phương sai số:
\[
\mathcal{L}(\theta) = \mathbb{E}_{(s,a,r,s') \sim \mathcal{D}}
\left[
\left(y - Q_{\theta}(s,a)\right)^2
\right]
\]
với $\mathcal{D}$ là bộ nhớ kinh nghiệm (\textit{bộ nhớ kinh nghiệm}).

\noindent\textbf{Kỹ thuật ổn định huấn luyện.}

\begin{itemize}
    \item \textbf{Phát lại kinh nghiệm}: lưu các mẫu chuyển trạng thái $(s,a,r,s')$ và lấy mẫu ngẫu nhiên theo mini-batch để giảm tương quan theo thời gian giữa các mẫu huấn luyện.
    \item \textbf{Mạng mục tiêu}: sử dụng mạng mục tiêu với tham số $\theta^-$, được cập nhật định kỳ từ mạng chính $\theta$, giúp giảm dao động trong quá trình học.
    \item \textbf{Khám phá epsilon-greedy}: tại bước huấn luyện, với xác suất $\epsilon$ tác nhân chọn hành động ngẫu nhiên để khám phá; ngược lại chọn hành động có Q lớn nhất để khai thác.
\end{itemize}

\section{Nghiên cứu liên quan}

\noindent Phần này tổng hợp các công trình liên quan đến tự phục hồi hệ thống, tối ưu điều khiển trên Kubernetes và ứng dụng học tăng cường trong vận hành hệ thống phân tán.
